\documentclass{ltxdoc}

\usepackage[T1]{fontenc}
\usepackage{array}
\usepackage{booktabs}
\usepackage{url}

\setlength{\emergencystretch}{2em}

\title{The \textsf{mention} Package}
\author{Ashwani Anand}
\date{2026/05/11}

\newcommand*\pkg[1]{\textsf{#1}}
\newcommand*\file[1]{\path{#1}}
\newcommand*\opt[1]{\texttt{#1}}
\newcommand*\csnamefmt[1]{\texttt{\symbol{92}#1}}

\begin{document}

\maketitle

\begin{abstract}
  The \pkg{mention} package adds targeted collaborative comments on top of
  \pkg{phfcc}.  Comments remain visible inline in the source document, while
  only the addressed collaborator sees margin highlights and a header summary
  of unresolved comments for the active viewer.
\end{abstract}

\tableofcontents

\section{Introduction}

Collaborative drafts often need comments that are visible in context but routed
to a specific co-author.  The \pkg{mention} package provides one command per
author, a recipient list for each comment, optional important-comment styling,
and a resolved flag that keeps old comments visible without counting them as
open work.

The package is intentionally small.  It depends on \pkg{phfcc} for the base
commenting environment and on common LaTeX packages for colors, margin notes,
and command parsing.  Automatic viewer detection is best-effort; projects that
need the most predictable behavior should use a local \file{mention-local.cfg}
file.

\section{Installation}

Place \file{mention.sty} where TeX can find it.  For a single project, putting
the file next to the main \file{.tex} document is enough:

\begin{verbatim}
project/
  paper.tex
  mention.sty
  mention-local.cfg.example
\end{verbatim}

Then load the package with:

\begin{verbatim}
\usepackage{mention}
\end{verbatim}

The package requires \pkg{phfcc}, \pkg{xparse}, \pkg{xcolor}, and
\pkg{marginnote}.  The optional automatic viewer detection path can use
\pkg{catchfile} and \pkg{texosquery} when they are available.

The built-in color cycle uses private copies of Paul Tol's muted
colorblind-friendly \pkg{colorblind} colors |T-Q-M1| through |T-Q-M9|.  If
\pkg{colorblind} has already defined those colors, \pkg{mention} copies them;
otherwise it uses the same RGB values internally.  The package does not load
\pkg{colorblind} and does not redefine document colors.

\section{Defining Authors}

\subsection{The author command factory}

Define one author command per collaborator:

\begin{verbatim}
\phfMakeTargetedCommentCommand[
  initials={AL},
  handle=alice,
  userid=alice-laptop
]{Alice}
\end{verbatim}

This defines a command named |\Alice|.  Its comments are printed inline using
the configured color.  If a comment targets the active viewer, a margin note is
also printed using the author's initials.

\begin{center}
\begin{tabular}{>{\ttfamily}ll}
\toprule
Option & Meaning \\
\midrule
initials & Short text printed in margin notes. \\
color & Any color expression accepted by \pkg{xcolor}. \\
handle & Stable collaborator handle used in recipient lists. \\
userid & Optional system user id mapped to the handle. \\
\bottomrule
\end{tabular}
\end{center}

If \opt{initials} is omitted, the command name is used.  If \opt{handle} is
omitted, the lower-cased command name is used.  If \opt{color} is omitted, the
package chooses the next color from its built-in colorblind-friendly cycle.

\section{Writing Comments}

Comments use the optional argument for recipient handles:

\begin{verbatim}
\Alice[bob]{Please check this lemma.}
\Alice[bob, claire]{Please verify this proof.}
\Alice![bob]{Please add the missing citation.}
\Alice[bob, resolved]{This has already been handled.}
\end{verbatim}

The form with |!| marks a comment as important.  Important comments use a
stronger filled margin label for the active viewer.

The special recipient-list item \opt{resolved} makes a comment dormant.  It
still appears inline, but it is not routed to anyone, does not get a margin
note, and does not contribute to unresolved counts.

Consecutive comments for the active viewer share one side-by-side margin note.
This keeps co-authors who comment at the same spot from covering each other in
the margin.

\section{Viewer Selection}

\subsection{Direct handle selection}

Set the active viewer directly by handle:

\begin{verbatim}
\phfSetViewer{bob}
\end{verbatim}

\subsection{User-id selection}

Set the active viewer by user id.  If the user id was registered through an
author's \opt{userid} option, it maps to that author's handle:

\begin{verbatim}
\phfSetViewerUser{bob-laptop}
\end{verbatim}

\subsection{Explicit user-id mappings}

Register or override an explicit user-id mapping:

\begin{verbatim}
\phfRegisterViewerUser{bob-laptop}{bob}
\end{verbatim}

Viewer selection is resolved in this order:

\begin{enumerate}
  \item an explicit |\phfSetViewer{...}| call;
  \item an explicit |\phfSetViewerUser{...}| call;
  \item a local \file{mention-local.cfg} file, if present;
  \item best-effort automatic detection.
\end{enumerate}

\subsection{Local Configuration}

For shared Git repositories, the recommended workflow is to keep
\file{mention-local.cfg} ignored and let each collaborator create their own copy
from \file{mention-local.cfg.example}.

\begin{verbatim}
cp mention-local.cfg.example mention-local.cfg
\end{verbatim}

A local config can select a handle:

\begin{verbatim}
\phfSetViewer{bob}
\end{verbatim}

or a user id:

\begin{verbatim}
\phfSetViewerUser{bob-laptop}
\end{verbatim}

The package loads \file{mention-local.cfg} automatically when the document
starts.

\subsection{Automatic Detection}

When no explicit viewer is set, \pkg{mention} tries to infer the compiling user.
With shell access, it first asks \file{kpsewhich} for common user variables:
|USER|, |USERNAME|, and |LOGNAME|.  If that does not work and \pkg{texosquery}
is available, it tries to derive a user id from the home-directory path.

This detection path is best-effort.  It may be unavailable on systems that
disable shell escape or do not ship the optional helper packages.  Manual
viewer selection through \file{mention-local.cfg} does not require shell escape.

\section{Header Summary}

\subsection{The standard header}

Enable a centered page-header summary:

\begin{verbatim}
\phfUseTargetedCommentHeader
\end{verbatim}

The summary reports unresolved comments for the active viewer, for example:

\begin{verbatim}
There are 2 unresolved comments for you.
\end{verbatim}

The count is written through the \file{.aux} file.  Compile twice after changing
comments or viewer selection.

\section{Complete Example}

\begin{verbatim}
\documentclass{article}
\usepackage[margin=1in]{geometry}
\setlength{\marginparwidth}{3cm}
\usepackage{xcolor}
\usepackage{mention}

\phfMakeTargetedCommentCommand[
  initials={AL},
  color=blue,
  handle=alice,
  userid=alice-laptop
]{Alice}

\phfMakeTargetedCommentCommand[
  initials={BO},
  color=orange,
  handle=bob,
  userid=bob-laptop
]{Bob}

\phfUseTargetedCommentHeader

\begin{document}

We should revise this proof
\Alice[bob]{Please check the induction step.}
before submission.

\Alice![bob]{Please add the missing citation.}
\Alice[bob, resolved]{This earlier point is done.}

\end{document}
\end{verbatim}

\section{Public Interface}

\begin{center}
\begin{tabular}{>{\ttfamily}ll}
\toprule
Command & Purpose \\
\midrule
\textbackslash phfMakeTargetedCommentCommand & define an author command \\
\textbackslash phfSetViewer & select the active viewer by handle \\
\textbackslash phfSetViewerUser & select the active viewer by user id \\
\textbackslash phfRegisterViewerUser & map a user id to a handle \\
\textbackslash phfUseTargetedCommentHeader & install the standard header \\
\textbackslash phfTargetedCommentSummary & typeset only the summary text \\
\bottomrule
\end{tabular}
\end{center}

The author-command factory has the form:

\begin{verbatim}
\phfMakeTargetedCommentCommand[options]{CommandName}
\end{verbatim}

The viewer commands have the forms:

\begin{verbatim}
\phfSetViewer{handle}
\phfSetViewerUser{userid}
\phfRegisterViewerUser{userid}{handle}
\end{verbatim}

\section{Limitations}

The response-marker commands
\begin{verbatim}
\phfMarkCommentResponded
\phfccDeclareResponded
\end{verbatim}
are retained only as compatibility stubs and emit a warning.  Use the
\opt{resolved} flag in a comment recipient list instead.

The package does not hide inline comments from non-recipients.  Targeting only
controls margin notes and unresolved counts.

\section{License}

Copyright \textcopyright{} 2026 Ashwani Anand.

This work may be distributed and/or modified under the conditions of the LaTeX
Project Public License, either version 1.3c of this license or any later
version.

This work has the LPPL maintenance status \texttt{maintained}.

The Current Maintainer of this work is Ashwani Anand.

This work consists of the files \file{README.md}, \file{mention.tex},
\file{mention.sty}, \file{mention-local.cfg.example}, and
\file{targeted-comments-demo.tex}, and the derived file \file{mention.pdf}.

\end{document}
